\documentclass[10pt]{article}

\usepackage{amsmath,amssymb,amsthm}

\newtheorem{definition}{Definition}
\newtheorem{problem}{Problem}
\newtheorem{theorem}{Theorem}
\newtheorem{remark}[theorem]{Remark}

\newcommand{\LQD}{\text{LQD}}
\newcommand{\LQDf}{\text{LQD}^{f}}
\newcommand{\LateQD}{\text{LateQD}}
\newcommand{\ALG}{\text{ALG}}
\newcommand{\OPT}{\text{OPT}}

\hyphenation{Niko-lenko}

\title{A Better Upper Bound for the Shared Memory Switch - Medium}
\author{Nikita Gaevoy}
\date{}

\begin{document}

\maketitle

\section{Preliminaries}

\subsection{The model}

A \emph{shared memory switch} has one input port and $N$ output ports. Each
output port $i\in\{1,\ldots,N\}$ has its own queue $Q^{(i)}$ of unbounded
individual size, subject to the single global constraint that the queues share
a common buffer of size $B$:
\[
\sum_{i=1}^{N}\bigl|Q^{(i)}\bigr|\ \le\ B .
\]
Packets are \emph{uniform}: they all have the same size and the same value, and
each packet is labelled with the index of the output port it is destined for.
Since packets are uniform, a queue is fully described by the number of packets
it holds, and we identify $Q^{(i)}$ with that number.

Packets arrive at discrete slotted time. Each time slot $t$ consists of two
phases.
\begin{enumerate}
    \item \textbf{Arrival phase.} A multiset $A_t$ of packets arrives; the
    input port may deliver arbitrarily many packets in one slot. Writing $B_t$
    for the set of packets in the buffer at the beginning of slot $t$ and
    $B_t'$ for the set after the arrival phase, a \emph{buffer management
    policy} is free to accept any subset of $A_t$ and to push out any subset of
    the buffer, subject to
    \[
    B_t'\subseteq B_t\cup A_t
    \qquad\text{and}\qquad
    |B_t'|\le B .
    \]
    \item \textbf{Transmission phase.} Every nonempty output queue transmits
    one packet, so that
    \[
    \bigl|Q^{(i)}\bigr|
    \ :=\
    \max\bigl\{0,\ \bigl|Q^{(i)}\bigr|-1\bigr\}
    \qquad\text{for every } i,
    \]
    and $B_{t+1}=\bigcup_{i=1}^{N}Q^{(i)}$.
\end{enumerate}

An \emph{instance} is a sequence $\tau:\mathbb{N}\to\mathbb{N}^{N}$, where
$\tau(t)_i$ is the number of packets arriving at slot $t$ destined for
$Q^{(i)}$. The instance has \emph{duration} $T(\tau)$ if no packets arrive
after slot $T(\tau)$. For an instance of finite duration and a policy $\ALG$,
we write $\ALG(\tau)$ for the total number of packets transmitted by $\ALG$ on
$\tau$; since at most $B$ packets remain in the buffer after the last arrival,
this is a finite quantity determined by the slots up to $T(\tau)+B$.

We denote by $\OPT$ the optimal offline (clairvoyant) policy, which knows the
entire instance in advance and has unbounded computational power. The goal is
throughput maximisation, so the competitive ratio is defined as

\begin{definition}[Competitive ratio]
For a policy $\ALG$,
\[
\alpha_{\ALG}
:=\sup_{\tau\,:\,T(\tau)<\infty}\alpha_{\ALG}(\tau),
\qquad\text{where}\qquad
\alpha_{\ALG}(\tau):=\frac{\OPT(\tau)}{\ALG(\tau)} .
\]
Thus $\alpha_{\ALG}\ge 1$, and smaller is better. A policy is
\emph{$\alpha$-competitive} if $\alpha_{\ALG}\le\alpha$.
\end{definition}

It is convenient to allow instances of infinite duration. If
$\tau_1\subseteq\tau_2\subseteq\cdots$ is an increasing sequence of instances,
with $\tau_n$ of length $n$ and the first $n$ slots of $\tau_{n+1}$ coinciding
with $\tau_n$, then $\tau=\bigcup_{n\ge 1}\tau_n$ is an \emph{infinite
instance} and we set
$\alpha_{\ALG}(\tau)=\lim_{n\to\infty}\alpha_{\ALG}(\tau_n)$. An infinite
instance may use an infinite number of output queues; if every queue receives
packets only during a finite window, such an instance can always be emulated
with finitely many queues by reusing queues that have been empty for more than
$B$ slots.

Throughout, $\alpha^{*}$ denotes the infimum of $\alpha_{\ALG}$ over all
deterministic online policies $\ALG$, i.e.\ the best competitive ratio
achievable online.

\subsection{The policies}

\begin{definition}[Longest Queue Drop]
The \emph{Longest Queue Drop} policy $\LQD$ accepts every arriving packet. If
this makes the total occupancy exceed $B$, then $\LQD$ repeatedly drops a
packet from a queue that is currently longest, breaking ties arbitrarily, until
the occupancy is again at most $B$.
\end{definition}

$\LQD$ is the natural online policy for this model: it equalises queue lengths,
and thereby keeps as many queues as possible nonempty and transmitting.

\begin{definition}[$\LateQD$]
Run the switch with an unlimited buffer, processing packets in LIFO order. In
that run each packet $p$ occupies the buffer over a time interval
$[b_p,e_p]$, where $e_p$ may be infinite. When the buffer overflows, the
clairvoyant policy $\LateQD$ pushes out a packet $p$ whose expected
transmission moment $e_p$ is latest; in particular packets with $e_p=\infty$
are pushed out first.
\end{definition}

\begin{theorem}[\cite{BDGN19}]\label{thm:lateqd}
$\LateQD$ is optimal, i.e.\ $\LateQD(\tau)=\OPT(\tau)$ for every instance
$\tau$. It is computable in linear time, and it remains optimal in the
fractional model.
\end{theorem}

Theorem~\ref{thm:lateqd} is what makes the offline side of a competitive ratio
effectively computable, and hence what makes hard instances searchable by
simulation.

\section{Problem formulation}

\begin{problem}[Upper bound]\label{prob:upper}
Prove that $\alpha^{*}<1.6918$, improving on the best upper bound known for
this problem~\cite{AEMV21,AEMV24}. Either of the following is a solution:
\begin{itemize}
    \item prove
    \[
    \alpha_{\LQD}\ \le\ c
    \]
    for an explicit constant $c<1.6918$, with no restriction on the number of
    output ports $N$ or on the buffer size $B$; or
    \item exhibit a deterministic online policy $\ALG$ with
    \[
    \alpha_{\ALG}\ <\ 1.6918 ,
    \]
    or, more ambitiously, with $\alpha_{\ALG}<\alpha_{\LQD}$ once the latter is
    known.
\end{itemize}
\end{problem}

\begin{remark}\label{rem:secondpolicy}
The first form solves the second, with $\ALG=\LQD$ itself. The second is worth
stating separately because $\LQD$ is at present the only policy for this
problem that is known to be competitive at all~\cite{AEMV24}, so even a second
competitive policy, with a worse constant, would be informative.
\end{remark}

\begin{remark}\label{rem:implies}
The most ambitious form of Problem~\ref{prob:upper}, a policy with
$\alpha_{\ALG}<\alpha_{\LQD}$, contradicts the complementary statement that
$\LQD$ is optimal among deterministic online policies,
$\alpha^{*}=\alpha_{\LQD}$; the other forms do not. Exactly one of that form
and that statement is true, and both are open.
\end{remark}

\begin{remark}\label{rem:uniform}
Bounds are asked for as absolute constants, uniform in $N$ and $B$. Bounds that
degrade to the trivial value $2$ as $N$ and $B$ grow are known, as are better
bounds for a fixed small number of output ports~\cite{KMO07,Matsakis15}.
Neither kind answers the problem above.
\end{remark}

\begin{thebibliography}{AEMV24}

\bibitem[AEMV21]{AEMV21}
Antonios Antoniadis, Matthias Englert, Nicolaos Matsakis, and Pavel Vesel{\'y}.
\newblock Breaking the barrier of 2 for the competitiveness of longest queue
  drop.
\newblock In {\em 48th International Colloquium on Automata, Languages, and
  Programming (ICALP 2021)}, volume 198 of {\em LIPIcs}, pages 17:1--17:20,
  2021.

\bibitem[AEMV24]{AEMV24}
Antonios Antoniadis, Matthias Englert, Nicolaos Matsakis, and Pavel Vesel{\'y}.
\newblock Breaking the barrier of 2 for the competitiveness of longest queue
  drop, 2024.
\newblock Full version of~\cite{AEMV21}; journal version in ACM Transactions on
  Algorithms, doi:10.1145/3676887.

\bibitem[BDGN19]{BDGN19}
Ivan Bochkov, Alex Davydow, Nikita Gaevoy, and Sergey~I. Nikolenko.
\newblock New competitiveness bounds for the shared memory switch, 2019.
\newblock arXiv:1907.04399.

\bibitem[KMO07]{KMO07}
Koji~M. Kobayashi, Shuichi Miyazaki, and Yasuo Okabe.
\newblock A tight bound on online buffer management for two-port shared-memory
  switches.
\newblock In {\em Proceedings of the 19th ACM Symposium on Parallelism in
  Algorithms and Architectures (SPAA)}, pages 358--364, 2007.

\bibitem[Mat15]{Matsakis15}
Nicolaos Matsakis.
\newblock {\em Approximation Algorithms for Packing and Buffering Problems}.
\newblock PhD thesis, University of Warwick, UK, 2015.

\end{thebibliography}

\end{document}
